% SGA 3，第 VII_B 篇，英文重建稿——第 1 轮。
% 范围：权威底本分篇第 37 页，第 1.5 节全文，止于 1.6 之前。
% 合并阅读本页码： 550.
% 再版印刷页码： 539.
% 本范围内原版页边定位： 510.
% 权威底本： Polo--Gille Exp7B-13oct24.pdf,
% SHA-256 D303690E0195654BBE7C3DAA5D89EA23E24CD709A3575217234AC19AE8836694.
% 曾参考 Jacob Reinhold 的 CC BY 4.0 Markdown 作为比对谱系。
\subsection*{1.5}\phantomsection \label{sga3:VIIB:1.5}\addcontentsline{toc}{subsection}{1.5. 平坦拓扑}

设 \(k\) 为伪紧环。

\medskip\noindent\textbf{定义。}\phantomsection \label{sga3:VIIB:1.5.definition} 若 \(\operatorname{Vaf}/k\) 中的态射族 \[
  f_i:X_i\longrightarrow X
\] 所诱导的态射 \[
  \coprod_i X_i\longrightarrow X
\] 满射且拓扑平坦，则称此族为\emph{满射拓扑平坦族}，这里该态射由各 \(f_i\) 诱导。这意味着每个 \(f_i\) 均拓扑平坦，且 \(X\) 的每个点都落在至少一个 \(X_i\) 的像中。

由 \SGARefLink{sga3:VIIB:1.3.3}{1.3.3} 可知，满射拓扑平坦族在 \(\operatorname{Vaf}/k\) 上定义预拓扑 \textup{（IV，\SGARefLink{sga3:iv:definition:4:2:5}{4.2.5}）}。相应拓扑称为 \(\operatorname{Vaf}/k\) 上的\emph{平坦拓扑}。

由 IV，\SGARefLink{sga3:iv:corollary:4:3:5}{4.3.5}，函子 \[
  F:(\operatorname{Vaf}/k)^{\mathrm{op}}\longrightarrow
  (\mathbf{Sets})
  % SGA3-VII-DIAG-094
\] 是平坦拓扑的层，当且仅当 \(F\) 将每个直和变为直积，且对每个满射拓扑平坦态射 \(f:X\to Y\)，序列\begin{center}\phantomsection \label{sga3:VIIB:1.5:flat-sheaf-sequence} % SGA3-VII-DIAG-094
% Authority: Exp7B-13oct24.pdf, local p. 37.
% Crop: figures/Exp7B_localp037_flat_sheaf_equalizer_sequence.png
% Locator: tex/components/30_expose_VIIB_section15_flat_topology_en.tex, Loop 1 line 48.
\(
F(Y)\xrightarrow{\ F(f)\ }F(X)
\mathrel{\substack{\xrightarrow{\ F(\operatorname{pr}_1)\ }\\[-.7ex]
                   \xrightarrow[\ F(\operatorname{pr}_2)\ ]{}}}
F(X\times_YX)
\)
\end{center}均正合。

由 IV，\SGARefLink{sga3:iv:subsection:4:5}{4.5}，\footnote{\label{sga3:VIIB:nde:69}下文由编者对原文作了修改。特别是将原来的陈述“若 \(d_0,d_1:X_1\rightrightarrows X\) 是等价关系，且 \(d_1\) 拓扑平坦，则商的形成与 \(h\) 交换”替换为下面的定理。}\SGARefLink{sga3:VIIB:1.3.1}{命题~1.3.1} 蕴含平坦拓扑比典范拓扑更粗。换言之，对 \(\operatorname{Vaf}/k\) 的每个对象 \(X\)，函子 \[
  h_X:T\longmapsto
  \operatorname{Hom}_{\operatorname{Vaf}/k}(T,X)
\] 是平坦拓扑的层。下文照例 \textup{（参见第一讲）} 将 \(X\) 与 \(h_X\) 等同。

由 IV，\SGARefLink{sga3:iv:lemma:4:6:5}{4.6.5}，\SGARefLink{sga3:VIIB:1.4.thm}{定理~1.4} 可以重新表述如下。

\medskip\noindent\textbf{定理。}\phantomsection \label{sga3:VIIB:1.5.theorem}\textit{设 \(k\) 为伪紧环，\(d_0,d_1:X_1\rightrightarrows X\) 为 \(\operatorname{Vaf}/k\) 中的等价对，\(X/X_1\) 为商形式簇 \textup{（即 \(\operatorname{Coker}(d_0,d_1)\)，参见 \SGARefLink{sga3:VIIB:1.2}{1.2}）}。若 \(d_1\) 拓扑平坦，则 \(X/X_1\) 表示平坦拓扑的商层。}
